\section{Introduction}
\label{sec:introduction}

The proliferation of Unmanned Aerial Vehicles (UAVs) in recent years has catalyzed their transition from niche applications to indispensable assets across a multitude of civil and military domains. These domains include, but are not limited to, precision agriculture, disaster management, infrastructure inspection, and persistent surveillance \cite{hafeez2023blockchain}. The operational efficacy of UAVs, particularly when deployed in cooperative swarms or as part of a Flying Ad-hoc Network (FANET), is critically dependent on the robustness, security, and efficiency of their underlying communication infrastructure. However, the very characteristics that make FANETs advantageous—such as their high mobility, dynamic topology, and reliance on wireless channels—also expose them to significant security vulnerabilities, including data tampering, eavesdropping, and single-point-of-failure risks associated with centralized control architectures \cite{luo2023escm}.

In response to these challenges, blockchain technology has emerged as a transformative paradigm with the potential to fundamentally enhance the security, autonomy, and trustworthiness of UAV communication systems. By leveraging its inherent properties of decentralization, immutability, and cryptographic security, blockchain offers a novel framework for establishing distributed trust, ensuring data integrity, and enabling secure, auditable transactions among UAVs without reliance on a central authority \cite{jagatheesaperumal2023blockchain}. The integration of blockchain into FANETs promises to mitigate prevalent security threats and foster a more resilient and cooperative operational environment.

This paper presents a comprehensive analysis of the application of blockchain technology to UAV communication networks. It critically examines the state-of-the-art, exploring various architectures, consensus mechanisms, and security protocols that have been proposed. Furthermore, this work introduces the **PrimeFusion FANET framework**, a novel methodology that synergistically combines a lightweight blockchain consensus protocol, Proof-of-Cooperation (PoC), with efficient data compression and a multi-layered security architecture. The objective of this research is to evaluate the viability of blockchain as a foundational technology for next-generation UAV networks, providing a detailed comparative analysis of its advantages and inherent limitations. Through this investigation, we aim to delineate the specific contexts in which blockchain integration provides tangible benefits and to identify the critical trade-offs that must be considered in practical deployments. The subsequent sections of this report will delve into the background of UAV communications and blockchain, articulate the core problem statement, review existing literature, detail the proposed PrimeFusion framework, and present a rigorous analysis of its performance and limitations before concluding with a discussion on future research directions.
